\section*{الفصل \ref{ch:g12:curly-arrows} --- آليات التفاعل: الأسهم المنحنية}

\begin{solution}{exo:g12:curly-arrows:1}
\ce{OH-}: مانح (أزواج حرة، شحنة سالبة)؛ ولا موقع مستقبِلًا فيه.
\ce{H+}: مستقبِل (لا \omterm{def:g9:inside-the-atom:nucleus}{إلكترونات} له إطلاقًا). \ce{CH2=CH2}: مانح (\omterm{def:g10:lewis-and-shape:multiple-bond}{الرابطة
الثنائية}). \ce{CH3-Cl}: مانح (أزواج \ce{Cl} الحرة)، ومستقبِل (الكربون،
$\delta^+$). \ce{H2O}: مانح (أزواج \ce{O} الحرة)؛ وذرتا الهيدروجين فيه،
وهما $\delta^+$، موقعان مستقبِلان ضعيفان.
\end{solution}

\begin{solution}{exo:g12:curly-arrows:2}
(a) \omterm{def:g12:curly-arrows:categories}{استبدال}؛ (b) إضافة؛ (c) حذف؛ (d) إضافة.
\end{solution}

\begin{solution}{exo:g12:curly-arrows:3}
يبدأ من زوج من \omterm{def:g9:inside-the-atom:nucleus}{الإلكترونات} (\omterm{def:g10:lewis-and-shape:lone-pair}{زوج حر} أو رابطة) وينتهي على
\omterm{def:g7:atoms-and-molecules:atom}{الذرة} التي يذهب إليها الزوج. ومن \omterm{def:g10:lewis-and-shape:lone-pair}{زوج حر} للأكسجين إلى كربون: يصير
الزوج رابطة \ce{O-C} جديدة.
\end{solution}

\begin{solution}{exo:g12:curly-arrows:4}
تعطي الأمونيا \omterm{def:g10:lewis-and-shape:lone-pair}{الزوج الحر} لنيتروجينها؛ والسهم يذهب من ذلك
\omterm{def:g10:lewis-and-shape:lone-pair}{الزوج الحر} إلى \ce{H+}. ويحمل \omterm{def:g7:chemical-reactions:reactant}{الناتج} \ce{NH4+} الشحنة $+1$.
\end{solution}

\begin{solution}{exo:g12:curly-arrows:5}
\omterm{def:g6:pure-substances-and-mixtures:species}{نوع كيميائي} يتكوّن في خطوة ويُستهلك في خطوة لاحقة؛ ولا يظهر في
المعادلة الإجمالية. وهو هنا الكربوكاتيون \ce{CH3-CH2+}.
\end{solution}

\begin{solution}{exo:g12:curly-arrows:6}
سهم من \omterm{def:g10:lewis-and-shape:lone-pair}{زوج حر} لأكسجين \ce{OH-} إلى الكربون؛ وسهم
من الرابطة \ce{C-Br} إلى البروم. الشحنات: $-1$ على اليسار
(\ce{OH-})؛ وعلى اليمين، الميثانول متعادل و\ce{Br-} يحمل $-1$:
فالمجموع $-1$ في الطرفين.
\end{solution}

\begin{solution}{exo:g12:curly-arrows:7}
البروم أكثر \omterm{def:g11:polarity-and-cohesion:electronegativity}{كهرسلبية} من الهيدروجين: فيأخذ الزوج
ويغادر أيونًا \ce{Br-} ($-1$)؛ ويرتبط الهيدروجين بكربون، ويبقى الكربون
الآخر بشحنة موجبة: \ce{CH3-CH2+} ($+1$).
\end{solution}

\begin{solution}{exo:g12:curly-arrows:8}
1.~\ce{CH3-CH2-OH + H+ -> CH3-CH2-OH2+}: المانح \omterm{def:g10:lewis-and-shape:lone-pair}{زوج حر} من
الأكسجين؛ والمستقبِل \ce{H+}.\par 2.~\ce{CH3-CH2-OH2+ -> CH3-CH2+ + H2O}: يذهب
زوج \ce{C-O} إلى الأكسجين الموجب (المستقبِل)، الذي يغادر به
على هيئة ماء.\par 3.~\ce{CH3-CH2+ -> CH2=CH2 + H+}: زوج رابطة \ce{C-H}
(المانح) ينتقل إلى الكربون الموجب (المستقبِل) ليكوّن \omterm{def:g10:lewis-and-shape:multiple-bond}{الرابطة الثنائية}،
ويغادر \ce{H+}. \omterm{def:g9:ions:ion}{وأيون} الهيدروجين المستعمل في الخطوة 1 يُعاد في الخطوة
3: فهو يظهر في الطرفين فيُحذف منهما.
\end{solution}

\begin{solution}{exo:g12:curly-arrows:9}
سهم من \omterm{def:g10:lewis-and-shape:lone-pair}{زوج حر} لأكسجين الماء إلى الكربون الموجب.
وبعد فقد \ce{H+}، يتكوّن الإيثانول \ce{CH3-CH2-OH}. والمعادلة الإجمالية:
\ce{CH2=CH2 + H2O -> CH3-CH2-OH}، وهي إضافة.
\end{solution}

\begin{solution}{exo:g12:curly-arrows:10}
يتحرك زوجان، في خطوة واحدة: تتكوّن الرابطة \ce{C-O}، وتنكسر الرابطة
\ce{C-Br}. ولا وسيط.
\end{solution}

\begin{solution}{exo:g12:curly-arrows:11}
الكلور أكثر \omterm{def:g11:polarity-and-cohesion:electronegativity}{كهرسلبية} ($3.16 - 2.55 = 0.61$ مقابل
$2.96 - 2.55 = 0.41$): فكربون الكلورو ميثان أشد
$\delta^+$، وهذا وحده كان سيجعله أسرع تفاعلًا. لكنه ليس كذلك: فالسرعة
تتوقف أيضًا على سهولة مغادرة \omterm{def:g10:electron-shells:family}{الهالوجين} بالزوج، والبروم
الأكبر يغادر بسهولة أكبر.
\end{solution}

\begin{solution}{exo:g12:curly-arrows:12}
إذا ارتبط الهيدروجين بالكربون 1: \ce{CH3-CH+-CH3}، الذي يعطي
2-كلورو بروبان \ce{CH3-CHCl-CH3}. وإذا ارتبط بالكربون 2:
\ce{+CH2-CH2-CH3}، الذي يعطي 1-كلورو بروبان \ce{CH2Cl-CH2-CH3}.
\end{solution}

\begin{solution}{exo:g12:curly-arrows:13}
السهم في الاتجاه الخطأ: فيجب أن يبدأ من المانح، \omterm{def:g10:lewis-and-shape:lone-pair}{زوج حر}
من \ce{OH-}، وأن ينتهي على الكربون المستقبِل. والرابطة الجديدة بين
الأكسجين والكربون، لا بين الأكسجين والبروم: فالبروم يغادر، آخذًا
زوج الرابطة \ce{C-Br}.
\end{solution}

\begin{solution}{exo:g12:curly-arrows:14}
كربون المجموعة \ce{C=O} في الحمض مرتبط بذرتي أكسجين
كهرسلبيتين: فهو $\delta^+$ بقوة، أي مستقبِل. وأكسجين
\omterm{def:g11:functional-groups:alcohol}{الكحول} يحمل $\delta^-$ وزوجين حرين: فهو مانح.
ويُفقد الماء إجمالًا؛ وتحل محلَّ مجموعة \ce{-OH} في الحمض مجموعةُ
\ce{-O-R} من \omterm{def:g11:functional-groups:alcohol}{الكحول}: إنه \omterm{def:g12:curly-arrows:categories}{استبدال}.
\end{solution}

\begin{solution}{exo:g12:curly-arrows:15}
تُدفع \omterm{def:g9:inside-the-atom:nucleus}{بإلكترونات} \omterm{def:g10:lewis-and-shape:multiple-bond}{الرابطة الثنائية}، فيتحرك زوج الرابطة
\ce{Br-Br} نحو البروم البعيد: فتصير \omterm{def:g7:atoms-and-molecules:atom}{ذرة} البروم القريبة
$\delta^+$، أي موقعًا مستقبِلًا. الخطوة الأولى: سهم من \omterm{def:g10:lewis-and-shape:multiple-bond}{الرابطة الثنائية}
إلى البروم القريب، وسهم من الرابطة \ce{Br-Br} إلى البروم
البعيد، الذي يغادر أيونًا \ce{Br-}؛ وتبقى شحنة موجبة على
الكربون الآخر \omterm{def:g10:lewis-and-shape:multiple-bond}{للرابطة الثنائية} السابقة.
\end{solution}

\begin{solution}{pb:g12:curly-arrows:1}
\textbf{1.} الإيثين: \ce{H2C=CH2}، \omterm{def:g10:lewis-and-shape:multiple-bond}{رابطة ثنائية} وأربع روابط \ce{C-H}.
بروميد الهيدروجين: \ce{H-Br} مع ثلاثة أزواج حرة على البروم.

\textbf{2.} \omterm{def:g10:lewis-and-shape:multiple-bond}{الرابطة الثنائية}: فزوجها الثاني منطقة غنية
\omterm{def:g9:inside-the-atom:nucleus}{بالإلكترونات}.

\textbf{3.} $2.96 - 2.20 = 0.76$: \omterm{def:g7:atoms-and-molecules:atom}{فذرة} \ce{H} تحمل $\delta^+$.

\textbf{4.} \omterm{def:g7:atoms-and-molecules:atom}{ذرة} الهيدروجين.

\textbf{5.} سهم من \omterm{def:g10:lewis-and-shape:multiple-bond}{الرابطة الثنائية} إلى هيدروجين \ce{HBr}؛
وسهم من الرابطة \ce{H-Br} إلى البروم.

\textbf{6.} الكربوكاتيون \ce{CH3-CH2+} ($+1$) و\ce{Br-} ($-1$).

\textbf{7.} قبلُ: $0 + 0 = 0$؛ وبعدُ: $+1 - 1 = 0$.

\textbf{8.} سهم من \omterm{def:g10:lewis-and-shape:lone-pair}{زوج حر} \omterm{def:g9:ions:ion}{لأيون} \ce{Br-} إلى الكربون
الموجب.

\textbf{9.} ستة (ثلاث روابط): ينقصه زوج ليكتمل ثُمانيّه، فهو يقبل
زوجًا من أي مانح قريب.

\textbf{10.} الكربوكاتيون \ce{CH3-CH2+}.

\textbf{11.} يتكوّن في الخطوة الأولى ويُستهلك في الثانية.

\textbf{12.} إضافة.

\textbf{13.} المجموعة فقط: فسلسلة الكربونين محفوظة،
\omterm{def:g10:lewis-and-shape:multiple-bond}{والرابطة الثنائية} تصير رابطة أحادية تحمل \ce{H} و\ce{Br}.

\textbf{14.} \qty{28.0}{g/mol}؛ \qty{108.9}{g/mol}.

\textbf{15.} $10.0 / 28.0 = \qty{0.357}{mol}$.

\textbf{16.} الإيثين: $0.357 - x$؛ \ce{HBr}: بفائض؛ البرومو إيثان: $x$.
فالإيثين هو المحِدّ: $x_{\max} = \qty{0.357}{mol}$.

\textbf{17.} $0.357 \times 108.9 = \qty{38.9}{g}$.

\textbf{18.} $0.75 \times 38.9 = \qty{29.2}{g}$.

\textbf{19.} نحو \qty{29}{g} من البرومو إيثان.
\end{solution}
